\documentclass[a4paper,12pt]{article}

% -------------------------------------------------
% Кодировка, язык, математика
% -------------------------------------------------
\usepackage[T2A]{fontenc}
\usepackage[utf8]{inputenc}
\usepackage[russian]{babel}
\usepackage{amsmath,amssymb,mathtools}
\usepackage{icomma}

% -------------------------------------------------
% Типографика и страница
% -------------------------------------------------
\usepackage{helvet}
\renewcommand{\familydefault}{\sfdefault}
\usepackage[a4paper,margin=2.15cm,headheight=15pt]{geometry}
\usepackage{microtype}
\usepackage{setspace}
\onehalfspacing
\usepackage{parskip}
\setlength{\parindent}{0pt}

% -------------------------------------------------
% Таблицы, списки
% -------------------------------------------------
\usepackage{tabularx,booktabs}
\usepackage{enumitem}

% -------------------------------------------------
% Графика и цвет
% -------------------------------------------------
\usepackage{tikz}
\usetikzlibrary{arrows.meta,calc}

\usepackage{xcolor}
\definecolor{accent}{HTML}{008080}
\definecolor{darkaccent}{HTML}{004D4D}
\definecolor{textgray}{HTML}{333333}
\definecolor{softgray}{HTML}{EEF3F3}
\definecolor{linegray}{HTML}{B9C8C8}

\color{textgray}

% -------------------------------------------------
% Заголовки
% -------------------------------------------------
\usepackage{titlesec}

\titleformat{\section}
  {\Large\bfseries\color{darkaccent}}
  {\thesection.}{0.65em}{}

\titleformat{\subsection}
  {\large\bfseries\color{accent}}
  {\thesubsection.}{0.6em}{}

\titlespacing*{\section}{0pt}{1.3em}{0.55em}
\titlespacing*{\subsection}{0pt}{1em}{0.4em}

% -------------------------------------------------
% Колонтитулы
% -------------------------------------------------
\usepackage{fancyhdr}
\pagestyle{fancy}
\fancyhf{}
\fancyhead[L]{\small Николь Саркисьянц}
\fancyhead[R]{\small Кредиты: смешанные схемы}
\fancyfoot[C]{\small\thepage}
\renewcommand{\headrulewidth}{0.4pt}
\renewcommand{\headrule}{%
  \hbox to\headwidth{\color{linegray}\leaders\hrule height \headrulewidth\hfill}}

% -------------------------------------------------
% Ссылки
% -------------------------------------------------
\usepackage{hyperref}
\hypersetup{
    colorlinks=true,
    linkcolor=darkaccent,
    urlcolor=accent
}

% -------------------------------------------------
% Собственные команды
% -------------------------------------------------
\newcommand{\important}[1]{\textbf{\color{darkaccent}#1}}
\newcommand{\rub}[1]{\,\text{#1}}
\newcommand{\mln}{\,\text{млн руб.}}

\begin{document}

% =================================================
% ТИТУЛЬНЫЙ ЛИСТ
% =================================================

\begin{titlepage}
\thispagestyle{empty}

\centering

\vspace*{3.2cm}

{\color{darkaccent}\rule{0.72\linewidth}{1.1pt}}

\vspace{1.2cm}

{\fontsize{27}{32}\selectfont\bfseries
Чек-лист\par}

\vspace{0.55cm}

{\fontsize{20}{25}\selectfont\bfseries
Кредиты: смешанные схемы погашения\par}

\vspace{0.35cm}

{\large
Таблица долга, платежи, переплата,\\
неравенства и отбор целых значений\par}

\vspace{1.2cm}

{\color{accent}\rule{0.42\linewidth}{0.8pt}}

\vspace{2.2cm}

{\large
\textbf{Лёвин Артём Александрович}\\[0.35em]
эксклюзивно для Николь Саркисьянц
\par}

\vspace{0.8cm}

{\normalsize
ЕГЭ по профильной математике
\par}

\vfill

{\large \today}

\vspace{1.6cm}

\begin{tikzpicture}
\draw[
    color=accent!55,
    line width=1.15pt
]
(0,0)
.. controls (3.0,0.75) and (8.5,-0.45) ..
(15.3,0.25);
\end{tikzpicture}

\vspace{0.7cm}

\end{titlepage}

% =================================================
% ОСНОВНОЙ ЧЕК-ЛИСТ
% =================================================

\section{Главная модель кредитной задачи}

В задачах со смешанной схемой график погашения задаётся индивидуально:
через таблицу, остатки долга, размеры отдельных платежей или дополнительные
условия.

Удобнее всего восстанавливать движение денег \important{по месяцам}.

Пусть

\[
S
\]

--- первоначальная сумма кредита,

\[
r\%
\]

--- процентная ставка за один период,

\[
p=\frac{r}{100}
\]

--- та же ставка в долях единицы.

Если в начале очередного месяца долг равен \(B_k\), то после начисления
процентов задолженность станет

\[
B_k(1+p).
\]

После платежа \(X_k\) останется

\[
B_{k+1}=B_k(1+p)-X_k.
\]

Отсюда основная формула для платежа:

\[
\boxed{X_k=B_k(1+p)-B_{k+1}}.
\]

Это центральное равенство для задач, в которых известны остатки долга.

\subsection{Стандартная структура таблицы}

\begin{table}[h]
\centering
\caption{Логика одной строки кредитной таблицы}
\renewcommand{\arraystretch}{1.35}
\begin{tabularx}{\linewidth}{>{\bfseries}l X}
\toprule
Элемент & Что записываем \\
\midrule
Начальный долг & \(B_k\) \\
Начисленные проценты & \(pB_k\) \\
Задолженность после процентов & \(B_k+pB_k=(1+p)B_k\) \\
Платёж & \(X_k\) \\
Остаток & \(B_{k+1}=(1+p)B_k-X_k\) \\
\bottomrule
\end{tabularx}
\end{table}

Ключевая связь между соседними строками:

\[
\boxed{\text{остаток после текущего платежа}
=
\text{долг в начале следующего периода}.}
\]

То есть

\[
B_{k+1}
\]

одновременно является остатком текущей строки и начальным долгом следующей.

% -------------------------------------------------

\section{Алгоритм решения}

Для большинства задач достаточно следующей последовательности.

\begin{enumerate}[leftmargin=1.8em,itemsep=0.45em]
    \item Определите первоначальный долг \(S\).
    \item Переведите процентную ставку:
    \[
    p=\frac{r}{100}.
    \]
    \item Запишите известные остатки долга \(B_2,B_3,\ldots\).
    \item Для каждого периода найдите начисленные проценты:
    \[
    pB_k.
    \]
    \item Найдите задолженность до платежа:
    \[
    (1+p)B_k.
    \]
    \item Восстановите платёж:
    \[
    X_k=(1+p)B_k-B_{k+1}.
    \]
    \item Для последнего периода используйте условие полного погашения:
    \[
    B_{n+1}=0.
    \]
    \item Составьте требуемое условием задачи уравнение или неравенство.
    \item Учтите целочисленность параметров и строгость неравенств.
\end{enumerate}

% -------------------------------------------------

\section{Последний платёж}

После последнего платежа кредит полностью погашен, поэтому

\[
B_{n+1}=0.
\]

Следовательно,

\[
X_n=(1+p)B_n.
\]

Последний платёж вполне может заметно отличаться от предыдущих.
График погашения определяется условием задачи.

% -------------------------------------------------

\section{Пример: остатки заданы в долях кредита}

Пусть кредит \(S\) выдан на шесть месяцев под \(5\%\) в месяц.

После последовательных платежей остаток составляет

\[
0{,}9S,\quad
0{,}8S,\quad
0{,}7S,\quad
0{,}6S,\quad
0{,}5S,\quad
0.
\]

Здесь

\[
p=0{,}05.
\]

Получаем:

\begin{table}[h]
\centering
\caption{Восстановление платежей}
\small
\renewcommand{\arraystretch}{1.3}
\begin{tabular}{ccccc}
\toprule
\(B_k\) &
\(0{,}05B_k\) &
\(1{,}05B_k\) &
\(X_k\) &
\(B_{k+1}\) \\
\midrule
\(S\) &
\(0{,}05S\) &
\(1{,}05S\) &
\(0{,}15S\) &
\(0{,}9S\) \\

\(0{,}9S\) &
\(0{,}045S\) &
\(0{,}945S\) &
\(0{,}145S\) &
\(0{,}8S\) \\

\(0{,}8S\) &
\(0{,}04S\) &
\(0{,}84S\) &
\(0{,}14S\) &
\(0{,}7S\) \\

\(0{,}7S\) &
\(0{,}035S\) &
\(0{,}735S\) &
\(0{,}135S\) &
\(0{,}6S\) \\

\(0{,}6S\) &
\(0{,}03S\) &
\(0{,}63S\) &
\(0{,}13S\) &
\(0{,}5S\) \\

\(0{,}5S\) &
\(0{,}025S\) &
\(0{,}525S\) &
\(0{,}525S\) &
\(0\) \\
\bottomrule
\end{tabular}
\end{table}

Сумма платежей:

\[
0{,}15S+0{,}145S+0{,}14S+
0{,}135S+0{,}13S+0{,}525S
=
1{,}225S.
\]

Переплата:

\[
1{,}225S-S=0{,}225S.
\]

Процент переплаты:

\[
\boxed{22{,}5\%}.
\]

% -------------------------------------------------

\section{Переплата}

Если все платежи равны

\[
X_1,X_2,\ldots,X_n,
\]

то общая сумма выплат составляет

\[
T=X_1+X_2+\ldots+X_n.
\]

Абсолютная переплата:

\[
\boxed{T-S}.
\]

Процент переплаты относительно первоначальной суммы кредита:

\[
\boxed{
\frac{T-S}{S}\cdot100\%.
}
\]

Полезная проверка:

\[
T-S
\]

равно сумме всех фактически начисленных процентов.

% -------------------------------------------------

\section{Неизвестная процентная ставка}

Иногда ставка \(r\) является искомой величиной.

Удобно сначала работать с

\[
p=\frac{r}{100},
\]

а переход к процентам выполнить после получения неравенства.

Рассмотрим кредит \(1\) млн руб. со следующими остатками:

\[
1;\quad 0{,}6;\quad0{,}4;\quad0{,}3;
\quad0{,}2;\quad0{,}1;\quad0.
\]

Платежи равны

\[
\begin{aligned}
X_1&=0{,}4+p,\\
X_2&=0{,}2+0{,}6p,\\
X_3&=0{,}1+0{,}4p,\\
X_4&=0{,}1+0{,}3p,\\
X_5&=0{,}1+0{,}2p,\\
X_6&=0{,}1+0{,}1p.
\end{aligned}
\]

Поэтому

\[
T=1+2{,}6p.
\]

Если требуется

\[
T<1{,}2,
\]

то

\[
1+2{,}6p<1{,}2,
\]

\[
2{,}6p<0{,}2,
\]

\[
p<\frac1{13}.
\]

Так как

\[
p=\frac r{100},
\]

имеем

\[
r<\frac{100}{13}
=
7\frac{9}{13}.
\]

При условии \(r\in\mathbb Z\) наибольшее допустимое значение:

\[
\boxed{r=7}.
\]

\subsection{Числовая прямая и отбор целого значения}

\begin{center}
\begin{tikzpicture}[x=1.25cm,y=1cm]
    \draw[-{Stealth[length=2.5mm]},line width=0.8pt]
        (0,0)--(8.8,0);

    \foreach \x in {0,1,...,8}{
        \draw (\x,0.08)--(\x,-0.08);
        \node[below=3pt] at (\x,0) {\small \x};
    }

    \draw[accent,line width=2pt] (0.15,0)--(7.67,0);

    \draw[accent,line width=1pt,fill=white]
        (7.69,0) circle (2.7pt);

    \node[above=5pt,text=darkaccent]
        at (7.69,0)
        {\(\displaystyle \frac{100}{13}=7\frac9{13}\)};

    \node[above=6pt,text=accent]
        at (7,0)
        {\(\boxed{7}\)};
\end{tikzpicture}
\end{center}

При выборе целого ответа используется положение числа на числовой прямой.
Обычное округление здесь не определяет ответ задачи.

% -------------------------------------------------

\section{Строгие и нестрогие условия}

Формулировки экономических задач требуют особого внимания.

\begin{table}[h]
\centering
\caption{Перевод словесного условия в неравенство}
\renewcommand{\arraystretch}{1.35}
\begin{tabular}{ll}
\toprule
Формулировка & Математическая запись \\
\midrule
\(A\) меньше \(B\) & \(A<B\) \\
\(A\) не больше \(B\) & \(A\le B\) \\
\(A\) больше \(B\) & \(A>B\) \\
\(A\) не меньше \(B\) & \(A\ge B\) \\
\bottomrule
\end{tabular}
\end{table}

Например, условия

\[
T<50
\]

и

\[
T\le50
\]

могут привести к разным целым ответам.

% -------------------------------------------------

\section{Условие «каждая выплата»}

Фраза

\[
\text{«каждая выплата больше 5 млн руб.»}
\]

означает одновременное выполнение всех ограничений.

Если после составления таблицы получены платежи

\[
0{,}55S,\qquad
0{,}475S,\qquad
0{,}5S,
\]

то требуется решить систему

\[
\begin{cases}
0{,}55S>5,\\
0{,}475S>5,\\
0{,}5S>5.
\end{cases}
\]

Из второго ограничения:

\[
S>
\frac{5}{0{,}475}
=
\frac{5000}{475}
=
\frac{200}{19}
=
10\frac{10}{19}.
\]

Именно это условие оказывается самым сильным.

Если \(S\) выражается целым числом миллионов рублей, то

\[
\boxed{S_{\min}=11}.
\]

% -------------------------------------------------

\section{Десятичные дроби и точный отбор ответа}

В задачах на кредиты дробь часто требуется сохранить в точном виде.

Например,

\[
\frac{100}{13}
=
7\frac9{13}.
\]

Такая запись сразу показывает:

\[
7<\frac{100}{13}<8.
\]

Для отбора наибольшего или наименьшего целого значения это надёжнее длинной
десятичной записи.

При выражениях вида

\[
\frac{5}{0{,}475}
\]

сначала можно устранить десятичную дробь:

\[
\frac{5}{0{,}475}
=
\frac{5000}{475}
=
\frac{200}{19}.
\]

% -------------------------------------------------

\section{Платежи, заданные через \(x\)}

В некоторых задачах размеры платежей задаются отношениями:

\[
x,\quad 3x,\quad 9x.
\]

Тогда общая сумма выплат равна

\[
x+3x+9x=13x.
\]

Если первоначальная сумма и ставка известны, значение \(x\) удобно находить
из условия полного погашения.

Для последней строки:

\[
\boxed{\text{задолженность перед последним платежом}
=
\text{последний платёж}.}
\]

Эквивалентная запись:

\[
B_{n+1}=0.
\]

% -------------------------------------------------

\section{Полезная свёрнутая формула}

Из рекуррентной формулы

\[
B_{k+1}=(1+p)B_k-X_k
\]

можно получить состояние долга после нескольких периодов.

Для трёх платежей:

\[
B_4
=
S(1+p)^3
-
X_1(1+p)^2
-
X_2(1+p)
-
X_3.
\]

При полном погашении

\[
B_4=0,
\]

поэтому

\[
\boxed{
S(1+p)^3
=
X_1(1+p)^2+
X_2(1+p)+X_3.
}
\]

Эта формула полезна, когда платежи выражены через одну неизвестную величину.

% -------------------------------------------------

\section{Типичные ошибки}

\begin{enumerate}[leftmargin=1.8em,itemsep=0.55em]
    \item Начисление процентов на первоначальную сумму \(S\) во всех строках.
    Проценты каждого периода считаются от текущего долга \(B_k\).

    \item Потеря связи
    \[
    B_{k+1}=\text{остаток после }k\text{-го платежа}.
    \]

    \item Вычитание величин в неправильном порядке.
    Платёж вычисляется по формуле
    \[
    X_k=(1+p)B_k-B_{k+1}.
    \]

    \item Использование \(r\) непосредственно как множителя.
    Для ставки \(5\%\)
    \[
    p=0{,}05.
    \]

    \item Пропуск условия
    \[
    B_{n+1}=0
    \]
    в последнем периоде.

    \item Неверная трактовка слов
    «меньше», «не больше», «больше», «не меньше».

    \item Округление промежуточной границы перед отбором целого ответа.

    \item Проверка только одного платежа при условии
    «каждая выплата».

    \item Поиск переплаты по отдельному месяцу.
    Сначала требуется сумма всех платежей.

    \item Пропуск ограничения на целочисленность \(S\) или \(r\).
\end{enumerate}

% -------------------------------------------------

\section{Краткий чек-лист перед ответом}

Перед сдачей задачи проверьте:

\begin{itemize}[leftmargin=1.7em,itemsep=0.4em]
    \item ставка переведена в долю:
    \[
    p=\frac r{100};
    \]

    \item проценты начислены на текущий долг;

    \item остаток одной строки перенесён в начало следующей;

    \item последний остаток равен нулю;

    \item платежи вычислены как
    \[
    X_k=(1+p)B_k-B_{k+1};
    \]

    \item условие задачи переведено в правильное уравнение,
    систему или неравенство;

    \item строгий или нестрогий знак выбран по тексту;

    \item целочисленность параметра учтена;

    \item итоговое значение удовлетворяет исходным ограничениям.
\end{itemize}

% =================================================
% ТРЕНИРОВОЧНЫЙ БЛОК
% =================================================

\newpage

\section*{Тренировочный блок}

Во всех задачах проценты начисляются перед очередным платежом.
Денежные величины, если специально не оговорено иное, указаны в миллионах
рублей.

\subsection*{Средний уровень}

\begin{enumerate}[leftmargin=2em,itemsep=1.15em]

\item
Кредит в размере \(2\) млн руб. выдан на три месяца под \(4\%\) в месяц.
После первого платежа долг равен \(1{,}6\) млн руб., после второго ---
\(1\) млн руб., после третьего кредит полностью погашен.

Найдите три платежа и их общую сумму.

\item
Кредит \(1\) млн руб. выдан на три месяца под \(5\%\) в месяц.
После первого платежа осталось \(0{,}8\) млн руб., после второго ---
\(0{,}5\) млн руб., после третьего долг равен нулю.

Найдите процент переплаты.

\item
Кредит \(S\) млн руб. выдан на три месяца под \(10\%\) в месяц.
После первого платежа долг составляет \(0{,}7S\), после второго ---
\(0{,}4S\), после третьего кредит погашен.

Каждый платёж должен быть больше \(4\) млн руб.
Известно, что \(S\) --- целое число.

Найдите наименьшее возможное значение \(S\).

\end{enumerate}

\subsection*{Выше среднего}

\begin{enumerate}[leftmargin=2em,itemsep=1.15em,start=4]

\item
Кредит \(1\) млн руб. выдан на пять месяцев.
После платежей долг последовательно составляет

\[
0{,}8;\quad0{,}6;\quad0{,}4;\quad0{,}2;\quad0.
\]

Месячная ставка равна \(r\%\), где \(r\) --- положительное целое число.

Найдите наибольшее \(r\), при котором общая сумма выплат
\[
\textbf{меньше }1{,}18\text{ млн руб.}
\]

\item
Условия предыдущей задачи сохранены, однако общая сумма выплат должна быть

\[
\textbf{не больше }1{,}18\text{ млн руб.}
\]

Найдите наибольшее целое значение \(r\).

\item
Кредит \(S\) млн руб. выдан на четыре месяца под \(10\%\) в месяц.
После платежей долг последовательно составляет

\[
0{,}75S,\quad0{,}5S,\quad0{,}25S,\quad0.
\]

Каждая выплата должна быть больше \(6\) млн руб.
Известно, что \(S\) --- целое число.

Найдите наименьшее возможное \(S\).

\item
Кредит \(245\) млн руб. выдан на три года под \(20\%\) годовых.
В конце каждого года совершается платёж.

Размеры платежей равны соответственно

\[
x,\quad2x,\quad4x.
\]

После третьего платежа кредит полностью погашен.

Найдите \(x\).

\item
Кредит \(S\) млн руб. выдан на три месяца под \(8\%\) в месяц.
После первого платежа долг равен \(0{,}7S\), после второго ---
\(0{,}4S\), после третьего кредит погашен.

Общая сумма выплат должна быть меньше \(40\) млн руб.
Известно, что \(S\) --- положительное целое число.

Найдите наибольшее возможное \(S\).

\item
Кредит \(2\) млн руб. выдан на три месяца под одинаковый процент \(r\)
в каждом месяце.

После первого платежа остаётся \(1{,}4\) млн руб.,
после второго --- \(0{,}8\) млн руб.,
после третьего кредит полностью погашен.

Общая сумма всех платежей равна \(2{,}3\) млн руб.

Найдите \(r\).

\end{enumerate}

\subsection*{Сложная задача}

\begin{enumerate}[leftmargin=2em,itemsep=1em,start=10]

\item
Кредит \(S\) млн руб. выдан на четыре месяца под \(r\%\) в месяц.
Известно, что

\[
S\in\mathbb Z,\qquad
r\in\mathbb Z,\qquad
1\le r\le20.
\]

После последовательных платежей долг составляет

\[
0{,}8S,\qquad
0{,}55S,\qquad
0{,}3S,\qquad
0.
\]

Каждая из четырёх выплат строго больше \(8\) млн руб., а общая сумма
выплат не превышает \(34\) млн руб.

Найдите наибольшее возможное значение \(r\).
Укажите также соответствующее целое значение \(S\).

\end{enumerate}

% =================================================
% ОТВЕТЫ
% =================================================

\newpage

\section*{Ответы}

\begin{table}[h]
\centering
\caption{Ответы к тренировочному блоку}
\renewcommand{\arraystretch}{1.5}
\begin{tabularx}{\linewidth}{c X}
\toprule
№ & Ответ \\
\midrule

1 &
\(
0{,}48;\;
0{,}664;\;
1{,}04
\) млн руб.;
общая сумма:
\(
2{,}184
\) млн руб. \\

2 &
\(
11{,}5\%
\) \\

3 &
\(
S=11
\) млн руб. \\

4 &
\(
r=5\%
\) \\

5 &
\(
r=6\%
\) \\

6 &
\(
S=22
\) млн руб. \\

7 &
\(
x=54
\) млн руб. \\

8 &
\(
S=34
\) млн руб. \\

9 &
\(
r=\dfrac{50}{7}\%
=
7\dfrac17\%
\) \\

10 &
\(
r=18\%,\qquad S=23
\) млн руб. \\

\bottomrule
\end{tabularx}
\end{table}

\end{document}